码：
身份证号、手机号码都是码。其中身份证号是一种系统码，可以解码成生日/户籍等信息，能找到一个错误因此dmin=2，手机号码不能纠错，d=1

循环码：
% 定义
循环码考虑的是循环不变性，借助生成多项式的原理进行表示。
% 定理9.1：C是循环码的充分必要条件是C是F[x]/(x^n-1)的理想
R_n = F[x]/(x^n-1) 是一个环
其实理论上来说你用什么多项式都能生成一种类型的码（多项式码），比如f(x) = x^n +1，这叫反循环码（Negacyclic Codes）。
乘 $x$ 相当于移位后，最后一位跑到第一位，但带上一个负号。

一般来说有以下定理：
F是域，F[x]是主理想整环，考虑F[x]/<f(x)>，则它是一个主理想环。并且生成元是f(x)的因式
这样循环码就是<g(x)>，g(x)是x^n-1的因式。被称为生成多项式
计算所有循环码，就是对x^n-1进行因式分解，找出所有的因式，然后每个因式生成一个循环码。

生成矩阵：很特殊的形状，
% 证明
这直接说明deg g(x) = n-k

c是一个码字当且仅当ch = 0, h是校验多项式。虽然校验多项式直接生成的码不是对偶码，但可以解方程知道其互反多项式是校验码的生成多项式。

分为系统和非系统两种编码：
非系统：直接用信息乘生成多项式得到码字
系统：先移位校验位个数量，然后对生成多项式取模，再减去。

不同线性码的本质是不同形状的线性子空间

布尔函数中的幂次重定义了：0^0 = 1, x^c就是看x和c是否相同

汉明码舍去了开头的1，本质上就是在说后面几位全部为零的这么一种码。很简单的一类线性空间

hamming码书上写的相当佶屈聱牙，能不能写一个好点的代替这个东西？
